\section{Analytic Strategy}
Because the frequency reports on culture talk are only available in the form of vague quantifiers we need a way to ``norm'' the individual's response to these items to hold constant idiosyncratic interpretations of the frequency categories.\footnote{There is an extensive literature on the use of vague quantifiers in survey data.  The main upshot of this literature is that structured heterogeneity in the interpretation of vague quantifiers does exist, and that this heterogeneity typically takes a simple form:  groups that are more likely to more frequently engage in the activity---in terms of numeric estimates---has a tougher standard for considering a given numeric frequency to constitute ``often'' or ``very often'' than groups that are less likely to engage in the activity.  In the context of arts-participation surveys, this means that the young and more educated probably interpret ``very often'' as implying a higher objective frequency of participation than older, less educated respondents, for instance.}  The aim is to determine whether persons in certain socio-demographic groups are more likely to report higher frequencies of culture talk than others, but the problem is that no ``objective'' estimate of frequency is available. My strategy is to use the self-reported frequencies on culture consumption as a baseline with which to determine whether persons engage in more culture talk than would be expected \emph{in relation to their own self-reported consumption frequencies in that cultural domain.} This normalizes the frequency reports not only in relation to any sort of within-respondent biases regarding the interpretation of such phrases as ``often'' and ``very often'' but also normalizes the response to whatever the baseline expectations regarding consumption is within cultural domains (e.g. movies versus art).  I can then examine the question of whether there exist any subgroup differences (e.g. young versus old, highly educated versus respondent who were not able to attain high levels of education) in the relative likelihood of engaging in cultural talk in relation to consumption as a way to ascertain whether there exists socio-demographic heterogeneity in the extent to which cultural talk is a normative part of each person's everyday life net of subgroup differences in absolute levels of engagement in the cultural form in question.

I proceed by arranging the data in the form of matched-pairs \citep[247-256]{agresti07}.  Each individual's response to the culture-talk item is matched against their own response to the consumption item (e.g. serving as their own ``control''), contributing two observations to the analysis.  An example of this arrangement is shown in table~\ref{tab:matched} for three hypothetical cases.   Accordingly, we can think of the (joint) outcome of interest as the subject-specific cross-classification of the response to the frequency of consumption item versus their response to the frequency of cultural talk item in that cultural domain, both expressed in terms of vague-quantifiers.  The data (collapsed across subjects in the form of a square four by four table) are shown in table~\ref{tab:contingency}.  The matched pairs design also has the advantage of providing a simple and parsimonious way of allowing me to test the hypothesis of whether the effect of different socio-demographic predictors is the same or different across the two outcomes \citep[Pp.]{agresti07}.

The model to be estimated is one where $logit[P(Y_{i1}\leq j)]=\alpha_{ij}$, and $logit[P(Y_{i2}\leq j)]=\alpha_{ij} + \beta$ \citep[254]{agresti07}, where $\beta$ is the subject-specific difference in the probability of engaging in talk over consumption:

\begin{equation}
logit[P(Y_{it}\leq j)]=\alpha_{ij} + \beta_t, t=1,2
\label{eq1}
\end{equation}•

Here I take advantage the ordered structure of the responses by restricting the comparisons to logits of cumulative probabilities (e.g. estimating a fixed-effects ordered logit model with proportional odds).  The dependent variable is the---identically coded---joint response to the culture consumption and culture talk items ($Y_{it}$) for individual.  Here $t=0$ for the observation corresponding to the consumption item and $t=1$ for the observation corresponding to the cultural talk item.  The coefficient estimate $\beta_t$ can be interpreted as giving us the difference in response probabilities between the talk and consumption items for each person (e.g. differences in the conditional marginal distribution of talk versus consumption):  \emph{that is the extent to which a person engages in more culture-talk relative to their consumption frequency or vice versa}.   If the parameter is positive ($\beta_t>0$) this indicates that culture-talk is more frequent in that cultural domain for that person in comparison to reported frequencies of consumption;  $\beta_t<0$ indicates the reverse: consumption is more prevalent than talk in that domain.  Finally if $\beta_t \approx 0$, then than means that both consumption and talk are equally prevalent. Note that the subject-specific ``fixed effects'' ($\alpha_{ij}$) condition on all idiosyncratic variation due to cross-individual differences ($i$) in the interpretation of the cutoff points ($j$).

If we are interested in cross-group differences in the likelihood of engaging in culture talk (relative to culture consumption), that is if we are interested in the question of whether \emph{certain groups use culture talk in that domain to a larger extent that we would expect given their reported consumption frequency}, we can just estimate equation~\ref{eq1} for different subgroups and then compare changes in the estimated magnitude of $\beta_t$.  Any observed positive trends can then be interpreted as indicating that culture-talk is more common within that subgroup than expected by chance, net of subject-specific variation in the relative likelihood of attaching different ``objective'' estimates to the subjective frequency indicators and relative to cross-group differences in the likelihood of participating more often in the cultural domain in question.  Thus, if we consider two groups of respondents ($i$ and $j$), if $\beta_t[i]>0$ and $\beta_t[j]<0$, then that means that members of group $i$, ``over-talk'' net of their consumption frequency, while members of group $j$ ``over-consume'' relative to their talk frequency in that domain.  If $\beta_t>0$  or $\beta_t<0$ for all groups (for instance if in some domain all respondents indicate higher frequencies of talk relative to consumption or the reverse), we may still compare relative differences across groups within cultural domains.  Thus, if $\beta_t[j]-\beta_t[i]>0$, we can conclude that members of group $j$ engage in more culture-talk relative to their own consumption frequency in comparison to members of group $j$, a vice versa if  $\beta_t[j]-\beta_t[i]<0$.  If $\beta_t[j]-\beta_t[i] \approx 0$, then the two groups engage in culture-talk relative to their own consumption at comparable rates within that domain.




